\section{Результаты}

Робот прошёл три этапа. Первый прогон, на ранней версии с хранилищем SQLite,
собрал 400 документов и выявил блокировку тогдашнего второго источника,
Кино-Театр.Ру. Демонстрационный прогон на MongoDB, до 100 документов, проверил
все механизмы на настоящих сайтах и измерил скорость. Наконец, полная обкачка
15--19 октября собрала итоговый корпус. Цифры корпуса ниже взяты из базы в её
нынешнем виде; их же показывают \texttt{make stats} и \texttt{make validate}.

\subsection*{Итоговый корпус}

\begin{table}[h]
\centering
\begin{tabular}{|l|r|r|r|}
\hline
\textbf{Показатель} & \textbf{Википедия} & \textbf{film.ru} & \textbf{Всего} \\
\hline
Документов                          & 64\,594 & 14\,972 & 79\,566 \\
Объём \enquote{сырых} данных        & 11,6 ГБ & 1,6 ГБ & 13,2 ГБ \\
Объём выделенного текста            & 672 МБ & 169 МБ & 841 МБ \\
Доля полезного текста               & 5,7\,\% & 10,5\,\% & 6,2\,\% \\
Средний размер документа            & 188 КБ & 110 КБ & 174 КБ \\
Средний объём текста, символов      & 6\,423 & 6\,530 & 6\,443 \\
Медиана объёма текста, символов     & 3\,629 & 5\,141 & --- \\
90-й перцентиль, символов           & 13\,404 & 11\,797 & --- \\
Максимум, символов                  & 202\,052 & 87\,623 & --- \\
\hline
\end{tabular}
\caption{Итоговый корпус}
\label{tab:corpus}
\end{table}

Требование в 30 тысяч документов перекрыто в два с половиной раза. Объём текста
считается в байтах UTF-8, доля --- как в первой лабораторной работе: байты
текста к байтам HTML.

Источники заметно различаются по форме распределения. У Википедии медиана почти
вдвое меньше среднего: большинство статей о фильмах короткие, а среднее тянут
вверх немногочисленные огромные статьи --- об известных людях (Дэвид Боуи, Пол
Маккартни, Аль Пачино) и франшизах, по 100--140 тысяч символов. Максимум, 200 тысяч,
--- у статьи \enquote{Россия}, одной из немногих посторонних (см. ниже). У film.ru статьи ровнее --- медиана близка к среднему, а самые
короткие (10-й перцентиль) почти вдвое длиннее, чем у Википедии. Минимум у обоих
источников --- около 1\,500 символов: это порог \texttt{min\_text\_length}.

Тематичность корпуса проверена по категориям, которые Википедия перечисляет в
HTML каждой статьи. У 494 статей (0,8\,\%) нет ни одной категории о кино. Большая
часть из них всё равно смежная с темой --- компьютерные игры и песни бондианы,
сказка о Пиноккио и персонажи франшиз, статьи о кинопроизводстве. Но попали и
совсем посторонние --- \enquote{Россия}, \enquote{Канада}, \enquote{Франция}: тематические категории
вида \enquote{Фильмы о\ldots} иногда включают и основную статью своей темы.

MongoDB хранит 14,0 ГБ данных в 4,5 ГБ на диске: движок WiredTiger по умолчанию
сжимает коллекции, и \enquote{сырой} HTML сжимается втрое.

\subsection*{Ход полной обкачки}

Полная обкачка заняла четыре дня календарного времени в семи запусках. Робота
останавливали вручную, чтобы обновить код, а ноутбук засыпал и один раз
перезагрузился; после каждого перерыва робот продолжал с того же места. Обкачка закончилась сама: очередь опустела.

Итоговое состояние очереди --- 125\,228 обработанных адресов. Из них около
25 тысяч --- страницы категорий Википедии и 545 --- страницы списков film.ru,
которые нужны только ради ссылок. Ещё 6\,388 адресов помечены \texttt{failed}: почти
все это несуществующие категории вроде \enquote{Фильмы Дании 1952 года}, на
которые ссылаются навигационные шаблоны-сетки \enquote{страна $\times$ год}. Они
проверены выборочно: Википедия отвечает \enquote{нет страницы с таким
названием}, статей в них нет. 554 адреса пропущены: 539 --- ссылки на правку
страниц, 9 запрещены \texttt{robots.txt}.

Средняя скорость в последнем, самом длинном запуске --- 0,51 страницы в секунду.
Паузы вежливости позволяют до полутора запросов в секунду на два источника;
остальное съедает ожидание ответа сервера. Кроме того, примерно каждый пятый
запрос уходит на страницу категории, которая документом не становится.

\subsection*{Что пошло не так}

Полная обкачка вскрыла несколько проблем, которых не было видно на сотне
документов. Каждая исправлена, и на каждую добавлен тест.

\textbf{Продолжение категорий закрыто robots.txt.} Страница категории
Википедии показывает только 200 статей, а ссылки на следующие страницы ведут в
\texttt{/w/index.php}, который запрещён роботам. Первая версия брала только
первые 200 статей каждой алфавитной категории, и очередь кончилась на 1\,400
документах. Решение --- обход всего дерева категорий о кино (раздел
\enquote{Устройство робота}).

\textbf{Уход из тематики через родительские категории.} Со страниц категорий фильмов
есть ссылки на \enquote{Категории произведений искусства по алфавиту} и
\enquote{Алфавитные категории, наполняемые шаблонами}, а через них --- на книги,
картины и вообще всё. Решение --- разрешать только категории со словами о кино в
названии.

\textbf{Блокировка Кино-Театр.Ру.} Сайт снова включил капчу --- после 549
нормально скачанных страниц, хотя пауза теперь была 5 секунд, а не одна. Значит,
защита считает общее число запросов, а не их частоту. Робот капчу распознал и в
базу не пустил, но продолжал запрашивать сайт каждые 5 секунд и списал как
недоступные 1\,606 адресов. После этого появилась пауза источника с нарастанием
(раздел \enquote{Устройство робота}). Блокировка не снималась больше суток, и за это
время с сайта удалось сохранить всего 53 документа. Источник заменён на
film.ru: он отдаёт обычный HTML, и за 15 тысяч запросов к нему не было ни одной
ошибки.

\textbf{Обрыв интернета.} Ночью несколько часов не работал DNS. Робот считал
каждую ошибку окончательной и за это время пометил \texttt{failed} 12\,320
исправных адресов. Их пришлось вернуть в очередь вручную. После этого временные
ошибки сети стали отдельным случаем с паузой и повторами, и дальше робот
переживал сбои сам: из 41 сетевого сбоя только два адреса не скачались за пять
попыток подряд, а после возврата в очередь скачались и они.

\textbf{Красные ссылки.} Навигационные шаблоны Википедии ссылаются на
несуществующие страницы, и такие ссылки ведут на страницу их создания
(\texttt{?action=edit\&redlink=1}). Робот обошёл около 6\,900 таких адресов
впустую, а одна служебная страница \enquote{Создание: \ldots} даже попала в базу
как документ. Решение --- отсекать адреса Википедии с параметрами; лишний документ
удалён.

\textbf{Двойной слэш в названии.} Нормализация схлопывала повторные слэши, и
статьи вроде \texttt{.hack//SIGN} получали несуществующий адрес. Повторные слэши
больше не трогаются.

\subsection*{Демонстрационный прогон и скорость}

Демонстрационный прогон на 100 документов (до замены второго источника) сохранил
их за 213 секунд при 170 скачанных страницах. На нём была найдена проблема
чередования источников: при строгом чередовании скорость была 0,39 страницы в
секунду, с выбором по готовности хоста --- 0,80, вдвое выше.

\subsection*{Первый прогон и блокировка Кино-Театр.Ру}

Первый прогон, на версии с SQLite, дал 400 документов за 933 обойдённых URL.
Из 119 адресов документов Кино-Театр.Ру сохранились лишь 17: остальные вернули
не документ, а страницу одинакового размера в 1041 символ:

\begin{verbatim}
Кино-Театр.Ру Роботам запрещено посещать наш сайт.
Пожалуйста, подтвердите, что вы не робот. Проверка: введите...
\end{verbatim}

Сайт закрыт сервисом ddos-guard, и после примерно 470 запросов защита перевела
робота на капчу. Хуже всего то, что робот подмены не заметил: ответ приходил с
кодом 200, и страница с капчой выглядела как любой другой документ. От попадания
в базу ста с лишним одинаковых заглушек спас случай --- порог
\texttt{min\_text\_length} в 1\,500 символов. Отсюда в роботе появился признак
страницы-заглушки \texttt{blocked\_marker} в конфиге источника. Окончательно
вопрос со вторым источником решила только замена на film.ru.

\pagebreak
